
\sourcenote{Genre: admissible rule (tier 2). The family machinery ---
the packaged node, the linkage, the two-case semantic fact, and its
conservativity warrant --- is stated in the \rulename{LiRPA} section's
family template and is not restated here.}

\begin{brecipe}[the \rulename{ReLU-Pos} instance]\label{rec:relupos-instance}
The node, and its intended rewrite over the family's sign-split fact:
\[
\inferrule*[Right=\rulename{ReLU-Pos}]{
  \mathcal{T} \rhd \Phi \vdash h \geq 0 \\
  \ReLU_{\{h\}\to\{y\}} \in H
}{
  \Phi \vdash y = h
}
\qquad\rightsquigarrow\qquad
\text{case analysis on }
(h \leq 0 \wedge y = 0) \vee (0 \leq h \wedge y = h).
\]
The \emph{active} branch discharges cleanly: its context carries the
conjunct $y = h$, which \rulename{Ax} and $\wedge$E extract --- the
goal verbatim. The \emph{boundary} branch does not: its context
carries $h \leq 0$ and $y = 0$, the premise supplies $h \geq 0$, and
the goal $y = h$ now needs the antisymmetry fact $h = 0$ --- which, as
Remark~\ref{rem:relupos-antisym} shows, no catalogue derivation
provides. The rewrite therefore does \emph{not} fully discharge into
catalogue nodes; like \rulename{PConstExact}, the rule is admissible
by reduction to verified machinery: the shipped warrant is the direct
lemma \texttt{sound\_reluPos}, in which the boundary case is a single
antisymmetry step in the model ($p \leq 0$ from the case, $0 \leq p$
from the premise, so $p = 0$ and the two cases agree).
\end{brecipe}

\begin{bprop}[semantic warrant]\label{prop:relupos-sound}
Suppose the linkage holds and the premise sequent --- same premise
formula, conclusion the atom $h \geq 0$ --- is semantically valid.
Then the sequent concluding $y = h$ is semantically valid.
\end{bprop}
\begin{proof}
Fix an admissible valuation $\sigma$ satisfying $\Phi$. The premise
eliminates entrywise: every entry $p$ of the preactivation tensor
satisfies $0 \leq p$. The linkage extraction
(\texttt{reluLinked\_facts}) gives the output tensor at the same
shape and, per entry, either $p \leq 0 \wedge y = 0$ or $0 \leq p
\wedge y = p$. In the second case $y = p$ is the goal entry. In the
first, $p \leq 0$ meets the premise's $0 \leq p$, so $p = 0$
(antisymmetry --- the model step the catalogue lacks) and $y = 0 = p$.
The conclusion atom is the entrywise equality of the two tensors at
their shared shape. That is \texttt{sound\_reluPos}, step for step.
\end{proof}

\begin{bremark}[the kink case has no catalogue derivation]
\label{rem:relupos-antisym}
The claim in Recipe~\ref{rec:relupos-instance} is an enumeration over
the catalogue's equality-capable conclusions. \rulename{Norm} and
\rulename{EqRefl} conclude only canonical-form identities ($h$ and $0$
are not canonically equal); \rulename{EqSym}, \rulename{EqElim}, and
\rulename{ArithAtom} at the $=$ comparator only \emph{transport}
equalities already proven; this family and the cell rules conclude
equalities only over their packaged applications' ports; the
combination rules (\rulename{Farkas}, \rulename{Syzygy},
\rulename{PBound}, \rulename{BernCone} with an order predicate)
conclude order atoms, never $=$; and \rulename{BernCone}'s degenerate
conic-predicate route would need a region atom pinning $h$ to the
single point $0$, which no rule derives --- region atoms denote
regions of positive-dimensional simplices. So $h \leq 0 \wedge h \geq
0 \vdash h = 0$ has no derivation, and the boundary branch cannot
discharge into nodes. This is not a defect to repair but a fact about
where the calculus keeps its arithmetic: order antisymmetry lives in
the verified model step, one line of \texttt{sound\_reluPos}, rather
than in a rule every certificate would pay wire format for.
\end{bremark}

\begin{bremark}[what the premise is for]
Without $\Phi \vdash h \geq 0$ the rule is false at any entry with $h
< 0$: there the node's semantics forces $y = 0 \neq h$. The premise is
consumed \emph{entrywise} --- a tensor-shaped preactivation must be
proven nonnegative in every entry, which is what a rank-$0$ constant
against a tensor atom asserts. No side condition excludes the boundary
$h = 0$, and none is needed: the two semantic cases agree there, which
is exactly the antisymmetry step above.
\end{bremark}

\begin{anchors}
Recipe \ref{rec:relupos-instance} warrant & \texttt{sound\_reluPos}\\
family linkage extraction & \texttt{reluLinked\_facts} (see the \rulename{LiRPA} section)\\
entrywise premise elimination & \texttt{QTensor.rel2\_const\_right\_elim}\\
the kink antisymmetry & the \texttt{le\_antisymm} step inside \texttt{sound\_reluPos}\\
payload structure & \texttt{PReLUPosData} (premise proof, two wires, linkage)\\
\end{anchors}
